\section{GIỚI THIỆU CHUNG}

Các mạng viễn thông hiện đại, đặc biệt trong bối cảnh 5G, trung tâm dữ liệu và mạng điều khiển bằng phần mềm, có quy mô lớn, nhiều lớp giao thức và phụ thuộc chặt chẽ giữa thiết bị, dịch vụ, cấu hình định tuyến, topology và chỉ số hiệu năng. Một sự cố nhỏ như sai ASN BGP, thiếu địa chỉ IP, link down, switch BMv2 dừng hoạt động hoặc rule SDN sai có thể gây mất kết nối diện rộng, suy giảm chất lượng dịch vụ và làm tăng chi phí vận hành. Trong thực tế, xử lý sự cố không chỉ là phát hiện bất thường, mà còn cần định vị thiết bị lỗi, xác định nguyên nhân gốc rễ (Root Cause Analysis -- RCA) và đưa ra hướng khắc phục có thể kiểm chứng. Các khảo sát về RCA tự động cho thấy bài toán này khó vì triệu chứng quan sát được thường gián tiếp, chịu nhiễu, và có thể bắt nguồn từ nhiều tầng cấu hình khác nhau~\citep{rca_survey2023}.

Các cách tiếp cận truyền thống thường dựa trên luật thủ công, kinh nghiệm chuyên gia hoặc mô hình học máy tĩnh. Những phương pháp này khó thích ứng khi topology thay đổi, khi xuất hiện dạng lỗi mới hoặc khi dữ liệu vận hành có nhiễu. Sự phát triển của mô hình ngôn ngữ lớn (LLM) và AI Agent mở ra khả năng mới: agent có thể đọc mô tả sự cố, sử dụng công cụ chẩn đoán, suy luận nhiều bước và giải thích kết quả bằng ngôn ngữ tự nhiên. Các nghiên cứu gần đây về LLM cho chẩn đoán mạng nhấn mạnh tiềm năng kết hợp tri thức miền, log vận hành và công cụ kiểm chứng thay vì chỉ dựa vào câu trả lời một lượt của mô hình~\citep{llm_network2024}. Tuy nhiên, agent thông thường vẫn mang tính tĩnh; nếu prompt, bộ nhớ hoặc mô tả công cụ không được cập nhật thì hiệu quả chẩn đoán dễ suy giảm khi môi trường mạng thay đổi.

Đề tài này nghiên cứu và xây dựng hệ thống \textbf{Self-Evolving AI Agents} cho xử lý sự cố mạng viễn thông trên nền tảng benchmark NIKA~\citep{nika2025}. Khác với agent tĩnh, self-evolving agent khai thác chính lịch sử xử lý nhiệm vụ, phản hồi và kết quả đánh giá làm nguồn dữ liệu để cập nhật bộ nhớ, công cụ hoặc chiến lược suy luận~\citep{selfevolving2024}. Mục tiêu của đề tài là biến mỗi lần chẩn đoán thành một nguồn kinh nghiệm giúp agent cải thiện ở các lần sau. Cụ thể, hệ thống cần tiếp nhận dữ liệu từ môi trường mạng giả lập, gọi các công cụ kiểm tra trạng thái mạng, suy luận nguyên nhân gốc rễ, gửi kết quả đánh giá và cập nhật tri thức chẩn đoán dài hạn.

Ý nghĩa thực tiễn của sản phẩm nằm ở việc giảm tải cho kỹ sư vận hành trong các bước lặp lại: kiểm tra kết nối, đọc cấu hình host/router, đối chiếu bảng định tuyến, phân tích log switch và tổng hợp bằng chứng. Trong bối cảnh vận hành thực tế, một trợ lý như vậy không thay thế chuyên gia, nhưng có thể đóng vai trò ``first responder'' để thu hẹp không gian giả thuyết, chuẩn hóa báo cáo sự cố và lưu lại kinh nghiệm đã được kiểm chứng cho các ca tương tự.

Đóng góp chính của sinh viên gồm:
\begin{itemize}
  \item Thiết kế quy trình agent chẩn đoán trên nền tảng NIKA/Kathará, tích hợp các dịch vụ MCP để agent có thể tương tác với router, host, switch P4/BMv2, telemetry và hệ thống chấm điểm.
  \item Đề xuất và hiện thực mô-đun \textbf{Procedural Memory}, lấy cảm hứng từ Skill-Pro, lưu trữ kỹ năng chẩn đoán dạng thủ tục và tái sử dụng kinh nghiệm phù hợp với ngữ cảnh sự cố.
  \item Đề xuất và hiện thực mô-đun \textbf{Tool Refinement}, lấy cảm hứng từ DRAFT, cho phép agent tự tinh chỉnh tài liệu mô tả công cụ MCP dựa trên các lần gọi công cụ thực tế mà không thay đổi mã nguồn công cụ.
  \item Thực nghiệm so sánh agent baseline và agent tự tiến hóa trên benchmark NIKA, đánh giá theo detection, localization, RCA, số bước, số lần gọi công cụ và chi phí token.
\end{itemize}
